\documentclass[11pt,a4paper]{article}
\usepackage[margin=28mm]{geometry}
\usepackage{amsmath,amssymb,amsthm,mathtools,booktabs}
\usepackage[T1]{fontenc}
\usepackage{lmodern}
\usepackage{hyperref}
\hypersetup{colorlinks=true,linkcolor=blue,citecolor=blue,urlcolor=blue,
pdftitle={Explicit uniform lower bounds for planar sets avoiding integral distances},
pdfauthor={lixiang90}}
\newtheorem{theorem}{Theorem}[section]
\newtheorem{lemma}[theorem]{Lemma}
\newtheorem{proposition}[theorem]{Proposition}
\newtheorem{corollary}[theorem]{Corollary}
\theoremstyle{remark}\newtheorem{remark}[theorem]{Remark}
\newcommand{\R}{\mathbb R}
\newcommand{\N}{\mathbb N}
\newcommand{\Z}{\mathbb Z}
\newcommand{\area}{\lvert}
\newcommand{\PhiR}{\sqrt R\left(\frac{\log\log R}{\log R}\right)^3}
\title{Explicit uniform lower bounds for planar sets avoiding integral distances\\
\large A trimmed and vertically compressed digit construction}
\author{lixiang90\thanks{Prepared with assistance from OpenAI Codex. The classical
digit method is attributed to S\'ark\"ozy. The existing upper-bound formalization
in the companion repository is by Allen Hart and is not a contribution of this paper.}}
\date{4 October 2026}
\begin{document}
\maketitle
\begin{abstract}
We construct a single explicit unbounded open subset $A$ of the Euclidean plane
with no positive integral distance between distinct points and prove that, for
every real $R\ge e$,
\[
 |A\cap B_R(0)|\ \ge\ \frac9{9604}\sqrt R
       \left(\frac{\log\log R}{\log R}\right)^3.
\]
The construction trims the alphabet of a digit product, reduces its vertical
scaling, and enlarges the surviving rectangles. Uniform estimates for entire
rectangles establish compatibility both within and between infinitely many
generations. A single rectangle of area $3/8$ suffices as the finite prefix.
Compared with the reference construction used here, the certified coefficient
improves by a factor $124416/2401$. Exact optimization within the specified
parameter family has asymptotic improvement factor $384e^{-3/2}$.
This is an explicit constant improvement with three logarithmic losses;
neither global optimality of the coefficient nor a lower bound of order
$\sqrt R$ is asserted.
\end{abstract}

\section{Problem, scope, and main result}
Write $B_R(0)=\{z\in\R^2:|z|<R\}$ and let
\[
 M(R)=\sup\{|E|:E\subset B_R(0)\text{ is measurable},\
       |x-y|\notin\Z_{>0}\ (x,y\in E,\ x\ne y)\}.
\]
Here $|E|$ denotes two-dimensional Lebesgue measure and all logarithms are
natural. Erd\H{o}s Problem 953 asks for the growth of $M(R)$ \cite{bloom}.
S\'ark\"ozy's digit construction supplies nearly square-root lower bounds
for robustly avoiding integral distances \cite{sarkozy1,sarkozy2};
Goenka and Moore give a modern treatment and higher-dimensional extensions
\cite{gm}. The square-root upper bound for the measurable problem is background
to the present work \cite{ulam}. We do not reprove or claim that upper bound.

\begin{theorem}\label{thm:main}
There is an explicitly specified unbounded open set $A\subset\R^2$ such that
no distance between distinct points of $A$ is a positive integer, and
\begin{equation}\label{eq:main}
 |A\cap B_R(0)|\ge\frac9{9604}\PhiR\qquad(R\ge e).
\end{equation}
Consequently $M(R)$ satisfies the same lower bound. The construction is fixed
independently of $R$.
\end{theorem}

The coefficient in \eqref{eq:main} is a guaranteed uniform coefficient, not a
proved optimum for $M(R)$ or even the largest coefficient available from this
set. The word ``optimal'' below refers exclusively to the finite choice of a
trim parameter within the stated family. The exponent $1/2$ is unchanged,
and the factor $(\log\log R/\log R)^3$ has not been removed.

\section{Definition of the fixed infinite construction}
For each integer $j\ge6$, put
\[
 B_j=2^j,\qquad C_j=B_j^2/16,
\]
let $n_j\ge2$ be the largest integer with $(2n_j)^{2n_j}\le C_j$, and set
$k_j=\lfloor C_j^{1/(2n_j)}\rfloor$. Thus
\begin{equation}\label{eq:params}
 2n_j\le k_j,\quad k_j^{2n_j}\le C_j<(k_j+1)^{2n_j},\quad
 C_j<[2(n_j+1)]^{2(n_j+1)}.
\end{equation}
Existence follows because $C_j\ge256$ and $(2n)^{2n}$ grows without bound.
For $j\ge14$ these parameters satisfy $k_j\ge8$.

Temporarily write $k=k_j$, $n=n_j$. For $1\le t\le\lfloor k/2\rfloor$, define
\begin{equation}\label{eq:height}
 q=k-t,\qquad H(k,t)=\min\left\{\frac1{256},\frac{4t^2}{9k^3}\right\}.
\end{equation}
For $\alpha=(\alpha_0,\ldots,\alpha_{n-1})\in\{0,\ldots,q-1\}^n$, put
\begin{align}
 x_{j,\alpha}&=B_j+\sum_{i<n}\alpha_i k^i,\label{eq:x}\\
 y_{j,\alpha}&=3B_j^2+k\sum_{i<n}\alpha_i k^{2i}.
 \label{eq:y}
\end{align}
Let $Q_{j,\alpha}$ be the open rectangle centred there, of full width $1/2$
and full height $H(k,t)$. Choose $t_j$ to maximize
\begin{equation}\label{eq:objective}
 D_j(t)=\frac12(k_j-t)^{n_j}H(k_j,t)
\end{equation}
over its finite integer range, taking the smaller $t$ in a tie. This rule is
an exact finite comparison of rational numbers.

Define the base rectangle
\[
 P=(-3/8,3/8)\times(-17/32,-1/32),\qquad |P|=3/8,
\]
and the complete set
\begin{equation}\label{eq:A}
 A=P\ \cup\ \bigcup_{j\ge14}\ \bigcup_{\alpha\in\{0,\ldots,k_j-t_j-1\}^{n_j}}
 Q_{j,\alpha}.
\end{equation}
This is a countable union of open sets, hence open and measurable. It is
unbounded since even the zero digit string has centre $(2^j,3\cdot4^j)$.
Compared with the reference tail, the alphabet loses $t-1$ symbols, the
vertical multiplier changes from $8k$ to $k$, and the rectangles are enlarged.
Thus \eqref{eq:A} is not obtained merely by adding points to the old tail.

\section{Entire rectangles within one generation}
We give estimates independent of $j,n$ and of any enumeration of pairs.
Take two different strings, let $r$ be their highest differing position, and
orient their difference so that its leading digit is $d\ge1$. Put
\[
 X=k^r,\qquad a=|\Delta y|,\qquad b=|\Delta x|,
 \qquad p=\frac{q-1}{k-1},\quad c=\frac{q-1}{k^2-1}.
\]
Both coordinate differences have the sign of the leading digit, since the
sum of all lower magnitudes is smaller than the leading place. In particular
$a,b$ are positive integers. Geometric sums give
\begin{equation}\label{eq:leading}
 (d-p)X\le b\le(d+p)X,\qquad
 k(d-c)X^2\le a\le k(d+c)X^2.
\end{equation}
Here $p<1$, $c<1/(k+1)\le1/9$, and $1\le d\le q-1\le k-2$.

\begin{lemma}\label{lem:lower}
For all these pairs, with $H=H(k,t)$,
\[
 (b-1/2)^2>2aH.
\]
\end{lemma}
\begin{proof}
If $r\ge1$, retaining the exact geometric remainder gives
\begin{align*}
 b&\ge dk^r-(k-t-1)\sum_{i<r}k^i\\
  &=(d-1)k^r+t\sum_{i<r}k^i+1
   \ \ge\ dt k^{r-1}+1.
\end{align*}
The last inequality uses $t\le k$ and
$\sum_{i<r}k^i\ge k^{r-1}$. With $L=tk^{r-1}>0$, \eqref{eq:leading} and
$H\le4t^2/(9k^3)$ give
\[
 2aH\le\frac89(d+1/9)L^2.
\]
Since $d\ge1$,
\[
 d^2-\frac89(d+1/9)\ge\frac1{81},
\]
so $(b-1/2)^2>2aH$. If $r=0$, then $a=kd$, $b=d$ and $2t\le k$;
thus $2aH\le2d/9$ whereas
$(d-1/2)^2-2d/9\ge d/36>0$.
\end{proof}

\begin{lemma}\label{lem:upper}
For every such pair,
\[
 (b+1/2)^2\le\frac{45}{32}a.
\]
\end{lemma}
\begin{proof}
Because $X\ge1$, \eqref{eq:leading} implies
$b+1/2\le(d+3/2)X$. The quadratic
$(d+3/2)^2-(k+2)d$ is convex in $d$ on $[1,k-2]$ and has equal endpoint
values $17/4-k<0$. Therefore
\[
 (b+1/2)^2\le(k+2)dX^2.
\]
Also $a\ge(8/9)kdX^2$, by $c\le1/9$ and $d\ge1$. Since $k\ge8$,
$(k+2)/k\le5/4$, proving the stated factor $45/32$.
\end{proof}

For points in these two rectangles, horizontal and vertical differences can
be written $b+u$, $a+v$ after taking absolute values, where $|u|\le1/2$,
$|v|\le H$. Lemma~\ref{lem:lower} shows
\[
 (b+u)^2+(a+v)^2\ge(b-1/2)^2+(a-H)^2>a^2.
\]
Lemma~\ref{lem:upper} and $H\le1/256$ show
\begin{align*}
 (b+u)^2+(a+v)^2
 &\le(b+1/2)^2+(a+H)^2\\
 &\le a^2+(45/32+1/128)a+H^2<(a+1)^2.
\end{align*}
Here $2-(45/32+1/128)=75/128>0$, $a\ge1$ and $H^2<1$.
All distances therefore lie in $(a,a+1)$, which contains no integer.
The diameter of one rectangle has square at most $1/4+(1/256)^2<1$.
Different horizontal centre coordinates are different integers, so the
rectangles are disjoint, and their exact combined area is \eqref{eq:objective}.

\section{Compatibility across infinitely many generations}
From geometric sums and \eqref{eq:params}, every centre in generation $j$ lies
in the frame
\begin{equation}\label{eq:frame}
 B_j\le x<5B_j/4,\qquad 3B_j^2\le y<7B_j^2/2.
\end{equation}
Indeed the horizontal increment is less than $k^n\le B_j/4$, and the vertical
increment is less than $k^{2n}\le B_j^2/16$; the displayed upper vertical
bound is a convenient looser one.

Consider frames of two generations $i<j$, so $B_i\le B_j/2$. If $a=y_j-y_i$
and $b=x_j-x_i$, then
\[
 b\ge3B_j/8\ge48,\qquad b^2\le a\le32b^2.
\]
For example, $a\ge17B_j^2/8$ and $b\le5B_j/4$ give the first square inequality;
$a\le7B_j^2/2$ and $b\ge3B_j/8$ give the second. Both $a$ and $b$ are integers.

Set $h=1/256$ and $\delta=1/512$. Two rectangles change $b$ by at most $1/2$
and $a$ by at most $h$. Consequently
\[
 \frac{95}{96}b\le b+u\le\frac{97}{96}b,\qquad |v|\le h.
\]
The lower comparison with $(a+\delta)^2$ is strictly positive since
\[
 (95/96)^2-64(h+\delta)=5569/9216>0,
 \qquad h^2-\delta^2>0.
\]
The upper comparison with $(a+1-\delta)^2$ is strictly positive since
\[
 2(1-h-\delta)-(97/96)^2=8915/9216>0,
 \qquad (1-\delta)^2-h^2>0.
\]
Thus all cross-generation distances lie in $(a+\delta,a+1-\delta)$.
This proves compatibility of every future pair of generations, without
truncating the infinite union.

For a centre $(b,a)$ in any generation $j\ge14$, \eqref{eq:frame} gives
\[
 b\ge64,\qquad a\le3b^2,\qquad 48b^2\le25a.
\]
The base rectangle is contained in
\[
 C_* =\{(x,y):x^2+(y+9/32)^2<1/4,\ -23/32<y<5/32\}.
\]
Comparing a point of $P\subset C_*$ to a new rectangle, the horizontal
perturbation is at most $3/4$ and the vertical perturbation belongs to
$[-81/512,369/512]$. Hence the horizontal square lies between
$(253/256)^2b^2$ and $(259/256)^2b^2$. The coefficients for the same two
square comparisons are respectively
\[
 (253/256)^2-6(81/512+\delta)=1033/65536>0,
\]
\[
 2(1-\delta-369/512)-(25/48)(259/256)^2
       =67871/3145728>0.
\]
The corresponding constant terms are positive. Distances from $P$ to every
new rectangle also lie in $(a+\delta,a+1-\delta)$. Finally the diameter
of $P$ has square $(3/4)^2+(1/2)^2=13/16<1$.
We have proved the distance-avoidance assertion of Theorem~\ref{thm:main}.

\section{A uniform gain in the area of every tail generation}
The reference area used to extract the previous uniform coefficient is
\[
 U_j=\frac{(k_j-1)^{n_j}}{108k_j^3}.
\]
It corresponds to full width $1/3$ and full height $1/(36k_j^3)$.
For $k\ge8$, the allowed choice $t=2$ has uncapped height $16/(9k^3)$.
Since our selected $t_j$ maximizes \eqref{eq:objective},
\begin{equation}\label{eq:gain}
 \frac{D_j(t_j)}{U_j}\ge96\left(\frac{k_j-2}{k_j-1}\right)^{n_j}
 \ge G:=96(6/7)^4=\frac{124416}{2401}.
\end{equation}
For completeness, put $f(x)=(x/2)\log((x-2)/(x-1))$. Using
$\log(1-u)\ge-u/(1-u)$ for $0<u<1$ gives, for $x\ge8$,
\[
 f'(x)=\tfrac12\log\frac{x-2}{x-1}
       +\frac{x}{2(x-2)(x-1)}
       \ge\frac1{2(x-2)(x-1)}>0.
\]
As $n\le k/2$ and the base of the power is below one,
\[
 \left(\frac{k-2}{k-1}\right)^n
 \ge\exp(f(k))\ge\exp(f(8))=(6/7)^4.
\]
This is a proof for every allowed parameter, not an extrapolation of
computed examples. The first tail generation has $j=14$, $(k,n,t)=(8,4,2)$
and exact area $D_{14}=9/4$; its reference area is $2401/55296$.

\section{Passage to every real radius}
Let $R_j=4^{j+1}$. The frame \eqref{eq:frame} and the rectangle dimensions
imply that every rectangle of generation $j$ lies strictly inside $B_{R_j}$.
For example $x+y+1<4B_j^2$ for $B_j\ge64$ bounds the Euclidean norm.
Every rectangle of generation $j+1$ lies outside $B_{R_j}$ since its lower
vertical coordinate exceeds $3\cdot4^{j+1}-1/512>R_j$.
In particular
\begin{equation}\label{eq:complete-area}
 |A\cap B_{R_j}|=3/8+\sum_{14\le\ell\le j}D_\ell\qquad(j\ge14).
\end{equation}

\begin{lemma}\label{lem:extraction}
Suppose a measurable distance-avoiding set $E$ satisfies
$|E\cap B_{R_j}|\ge G U_j$ for every $j\ge6$, with $G>0$. Then
\[
 |E\cap B_R|\ge\frac{G}{55296}\PhiR\qquad(R\ge16384).
\]
\end{lemma}
\begin{proof}
Choose the unique $j\ge6$ with $R_j\le R<R_{j+1}=4R_j$.
Bernoulli's inequality, $2n\le k$, and $C_j<(k+1)^{2n}$ give
\[
 (k-1)^n\ge k^n/2\ge(k+1)^n/4>\sqrt{C_j}/4=B_j/16.
\]
Therefore
\begin{equation}\label{eq:U-radius}
 U_j>\frac{B_j}{1728k_j^3}>\frac{\sqrt R}{6912k_j^3}.
\end{equation}
For $n_j\ge6$ we have the explicit base bound
\begin{equation}\label{eq:klog}
 k_j<\frac{2\log R}{\log\log R}.
\end{equation}
Here are details of the bound and its finite cases. Write $T=\log R$.
The budget $2n\log(2n)<T$ implies $n<3T/(4\log T)$: for
$s=3T/(2\log T)$, the tangent inequality for $\log$ gives
$\log s>(2/3)\log T$, hence $s\log s>T$, and monotonicity of $x\log x$
completes the inversion. For $n\ge16$, induction gives
$2(n+1)\le(5/4)^n$, beginning with $34\cdot4^{16}<5^{16}$.
Together with \eqref{eq:params}, this yields
$k<(5/2)(n+1)\le(8/3)n<2T/\log T$.
For $6\le n\le15$, the following integer maxima $K$ bound $k$:
\[
\begin{array}{c|rrrrrrrrrr}
n&6&7&8&9&10&11&12&13&14&15\\\hline
K&21&23&25&27&29&32&34&36&38&40
\end{array}
\]
They follow by the exact inequality
$[2(n+1)]^{2(n+1)}\le(K+1)^{2n}$.
Set $T_0(k)=\log10+2n\log k$. The function
$k\log T_0(k)/T_0(k)$ is increasing for $k\ge2n$, while
$\log T/T$ is decreasing for $T\ge T_0(k)>e$.
For each displayed $K$, the inequality
$K\log T_0(K)<2T_0(K)$ is proved by rational logarithm bounds;
these ten inequalities are also present as kernel-checked lemmas
in the companion source. Their premises apply here because
$R\ge64C_j\ge10k_j^{2n_j}$.
Substituting \eqref{eq:klog} into \eqref{eq:U-radius} proves the lemma for $n_j\ge6$.

The remaining generations are exactly $6\le j\le23$. They can be checked
with the short table in Appendix~\ref{app:early}: throughout each whole
interval $[R_j,R_{j+1})$, use
$\sqrt R<2^{j+2}$ and the indicated upper bound $q_0$ for
$\log\log R/\log R$. Every table row satisfies the exact rational inequality
\[
 55296 U_j>2^{j+2}q_0^3.
\]
These are interval checks and so cover real as well as integer radii.
Finally area is monotone with radius, giving the claim.
\end{proof}

For $j\ge14$, \eqref{eq:gain} proves the premise of
Lemma~\ref{lem:extraction}. For $6\le j\le13$, the base rectangle alone
suffices: the exact values of $GU_j$ are
\[
\begin{array}{c|rrrr}
j&6&7&8&9\\\hline
GU_j&162/2401&18432/300125&9/196&115200/3195731\\[2pt]
j&10&11&12&13\\\hline
GU_j&2000/7203&9/28&104976/300125&2662/7203
\end{array}
\]
Each is strictly below $3/8$. Thus for $R\ge16384$ we get the coefficient
$G/55296=9/9604$. For $e\le R<16384$, the entire base rectangle lies in
$B_R$, while $\sqrt R\le128$ and $0\le\log\log R/\log R\le1$.
Since $128\cdot9/9604<3/8$, this covers the remaining radii and completes
the proof of Theorem~\ref{thm:main}.

The extremal quantity $M(R)$ independently satisfies $M(R)=\pi R^2$ for
$0<R\le1/2$ and $M(R)\ge\pi/4$ for $R\ge1/2$, by open disks of radius
at most $1/2$. Combining these with \eqref{eq:main} gives a lower bound for
$M(R)$ at every positive radius. The witnesses for those small-radius
bounds need not coincide with the fixed infinite set $A$.

\section{Optimization within the certified family}\label{sec:optimality}
The guaranteed coefficient used only the admissible choice $t=2$.
Optimization generally does better. From \eqref{eq:params},
\[
 2n\le k<[2(n+1)]^{1+1/n},\qquad k/(2n)\longrightarrow1.
\]
In the uncapped region, maximizing \eqref{eq:objective} is equivalent to
maximizing $t^2(k-t)^n$. Its logarithmic derivative is
$2/t-n/(k-t)$, strictly decreasing, with unique zero
\[
 t_{\rm real}=\frac{2k}{n+2}\longrightarrow4.
\]
The height cap begins at $t=\sqrt{9k^3/1024}$, which tends to infinity.
Above that point the capped area decreases with $t$. Thus the cap does not
affect the eventual maximizer. Comparison of the finitely many adjacent
integers, or their limiting objective ratios, shows that the unique optimal
integer is eventually $t=4$. For clarity, the limiting objective is
$t^2e^{-t/2}$, whose unique real maximum is $t=4$; its neighbouring integer
values at $3$ and $5$ are strictly smaller, and log concavity excludes all
other integers.

Consequently, for the reference per-generation areas,
\[
 \frac{D_j(t_j)}{U_j}
 =24t_j^2\left(\frac{k_j-t_j}{k_j-1}\right)^{n_j}
 \longrightarrow384e^{-3/2}=85.68198149\ldots.
\]
The fraction of surviving digit strings tends to $e^{-3/2}$, approximately
$22.313\%$, while the area of each surviving rectangle is multiplied by
$384$. The same ratio holds for cumulative areas at completed-generation
radii, since a finite prefix is negligible and the reference cumulative
area diverges. No ratio limit at arbitrary partially completed radii is
claimed.

This family cannot remove a logarithmic loss. The elementary parameter bound
$k<9n$ follows from \eqref{eq:params}: for $n\ge2$,
$[2(n+1)]^{1+1/n}<(9n)$ (one may use
$2(n+1)\le3^n$ and $6(n+1)\le9n$).
For any allowed $t$,
\[
 \frac{D_j(t)}{U_j}\le24t^2
    \left(1-\frac{t-1}{k-1}\right)^n
 \le24t^2 e^{-(t-1)/9}.
\]
The right side is bounded independently of $j$ and $t$. Thus arbitrary
choices of the trim parameter within this uniform-height digit product
can improve the multiplier only by a bounded amount. Non-product digit
selection or variable shapes would require a new argument.

\section{Verification, reproducibility, and attribution}
The companion repository is
\par\noindent\url{https://github.com/lixiang90/erdos953-lower-bound-lean}.
\par\noindent
Its formal proofs use Lean~4.32.2 and Mathlib commit
\par\noindent\texttt{905b95818eb32af7874a58b427f50c1711a5e96c}.
\par\noindent
The uniform inequality for the single fixed set, for every real $R\ge e$,
has been checked in Lean as
\texttt{Erdos953Lower.retreatInfinity\_uniform\_lower}.
The combined existence statement, including openness, unboundedness,
and exclusion of every positive integral distance, is
\texttt{Erdos953Lower.fixed\_unbounded\_uniform\_lower}.
Its extremal-area consequence is
\texttt{Erdos953Lower.retreat\_lower\_Mopen}.
The supporting modules verify parameter existence, the finite maximizing
trim rule and its guaranteed gain, the area of each generation, exclusion
of all positive integral distances within and between generations, and
compatibility with the base rectangle. The checked theorems use only
Lean's standard axioms \texttt{propext}, \texttt{Classical.choice}, and
\texttt{Quot.sound}; there are no proof placeholders.
The asymptotic maximizing trim $t=4$ and the limiting improvement factor
$384e^{-3/2}$ in Section~\ref{sec:optimality} are proved analytically in
this paper and are not part of the Lean verification claim.
A rational certificate audit, finite pair enumeration, and a Lean theorem
are distinct forms of evidence. Run \texttt{lake build Erdos953Retreat}
to check this lower construction independently of the upper-bound project.
The companion \texttt{scripts/build\_paper.py} runs LaTeX until citations
and cross-references are resolved and rejects unresolved-reference warnings.

The enlarged finite prefix from previous computations can be appended in
place of $P$. It consists of the clipped disk $C_*$ and 7,813 rational
rectangles. Its compatibility with the new tail is certified by rational
envelopes covering every future source frame. Those data give stronger
finite-radius area values, but are unnecessary for the main construction
\eqref{eq:A} and its coefficient. In particular the present analytic proof
does not rely on enumerating those 7,813 rectangles.

The classical digit principle is due to S\'ark\"ozy; our changes are the
trimmed alphabet, compressed vertical scaling, whole-rectangle estimates,
infinite compatible assembly, and explicit constant extraction.
The separately authored Lean upper-bound proof by Allen Hart is fetched
at commit \texttt{0d031f7212150df788f4f}\allowbreak\texttt{a38c26cfc3fc729f3d0} for the
earlier growth-exponent theorem. It is neither used in this paper's lower
construction nor redistributed as our source. Award eligibility and a
full extremal solution are not conclusions of this paper.

\appendix
\section{Short exact checks for the initial scales}\label{app:early}
The following table contains the entire finite part of the complete-generation
radius check. Here $q_0$ is an upper bound valid throughout the real interval
$[4^{j+1},4^{j+2})$. All rows satisfy
$512(k-1)^n>2^{j+2}k^3q_0^3$, equivalently the inequality used in
Lemma~\ref{lem:extraction}.
\begin{center}
\begin{tabular}{rrrr|rrrr|rrrr}
\toprule
$j$&$k$&$n$&$q_0$&$j$&$k$&$n$&$q_0$&$j$&$k$&$n$&$q_0$\\
\midrule
6&4&2&$1/4$&12&10&3&$1/4$&18&16&4&$1/6$\\
7&5&2&$1/4$&13&12&3&$1/6$&19&10&5&$1/8$\\
8&8&2&$1/4$&14&8&4&$1/6$&20&12&5&$1/8$\\
9&11&2&$1/4$&15&9&4&$1/6$&21&13&5&$1/8$\\
10&6&3&$1/4$&16&11&4&$1/6$&22&16&5&$1/8$\\
11&8&3&$1/4$&17&13&4&$1/6$&23&18&5&$1/8$\\
\bottomrule
\end{tabular}
\end{center}
To verify the bounds $q_0$, use
\[
 \log T\le T/4\ (T\ge9),\qquad
 \log T\le T/6\ (T\ge18),\qquad
 \log T\le T/8\ (T\ge27).
\]
They follow from the tangent inequality
$\log T\le\log T_0+(T-T_0)/T_0$ at $T_0=9,18,27$, respectively,
and the rationally verifiable bounds $\log3<9/8$, $\log18<3$.
Also $\log2>27/40$, so the respective first-generation log radii
$14\log2$, $28\log2$, $40\log2$ exceed $9,18,27$.
For a positive rational $x$, scale to $x=2^m y$, $1\le y\le2$, and use
$z=(y-1)/(y+1)$ in
\[
 2\sum_{r=0}^{s-1}\frac{z^{2r+1}}{2r+1}
 \le\log y\le
 2\sum_{r=0}^{s-1}\frac{z^{2r+1}}{2r+1}
 +\frac{2z^{2s+1}}{(2s+1)(1-z^2)}.
\]
Every comparison then reduces to exact rational arithmetic. The companion
Lean logarithm certificate uses a slightly looser remainder and six terms,
which are sufficient for the constants displayed here.

\clearpage
\begin{thebibliography}{9}
\bibitem{bloom} T.~F.~Bloom, \emph{Erd\H{o}s Problem 953},
\url{https://www.erdosproblems.com/953}, accessed 4 October 2026.
\bibitem{sarkozy1} A.~S\'ark\"ozy, \emph{On distances near integers. I},
Studia Sci. Math. Hungar. \textbf{11} (1976), 37--50.
\bibitem{sarkozy2} A.~S\'ark\"ozy, \emph{On distances near integers. II},
Studia Sci. Math. Hungar. \textbf{11} (1976), 105--111.
\bibitem{gm} R.~Goenka and K.~Moore,
\emph{Point sets avoiding near-integer distances}, arXiv:2605.06621v1 (2026),
\url{https://arxiv.org/abs/2605.06621}.
\bibitem{ulam} \emph{The order of growth of planar sets avoiding integer distances},
ULAM research preprint (2026),
\url{https://www.ulam.ai/research/erdos953-short.pdf}.
\bibitem{hart} A.~Hart, \emph{Erd\H{o}s 953 formalization},
\url{https://github.com/AllenGrahamHart/FormalConjectures-Bench/tree/0d031f7212150df788f4fa38c26cfc3fc729f3d0/formalizations/erdos953}.
\end{thebibliography}
\end{document}
